\section*{第 \ref{ch:g1:floating-and-sinking} 章——漂浮与下沉}

\begin{solution}{exo:g1:floating-and-sinking:1}
\omterm{def:g1:floating-and-sinking:words}{漂浮}的：软木塞、干树枝、塑料鸭。\omterm{def:g1:floating-and-sinking:words}{下沉}的：
钥匙、弹珠、石头。
\end{solution}

\begin{solution}{exo:g1:floating-and-sinking:2}
溅起的水花甚至可能把\omterm{def:g1:floating-and-sinking:words}{漂浮}的东西短暂压入水下，
所以扔进水里的物体可能会骗人。轻轻放下并
等待，才能让测试公平。
\end{solution}

\begin{solution}{exo:g1:floating-and-sinking:3}
金属勺\omterm{def:g1:floating-and-sinking:words}{下沉}，塑料勺\omterm{def:g1:floating-and-sinking:words}{漂浮}。诺拉应该看一个物体
由什么\emph{材料}制成，而不是它的名字或用途。
\end{solution}

\begin{solution}{exo:g1:floating-and-sinking:4}
重且\omterm{def:g1:floating-and-sinking:words}{漂浮}的：树干（或船）。轻且\omterm{def:g1:floating-and-sinking:words}{下沉}的：
硬币（或钥匙）。所以“重的东西会\omterm{def:g1:floating-and-sinking:words}{下沉}”是错误的规则——
仅凭重力无法决定沉浮。
\end{solution}

\begin{solution}{exo:g1:floating-and-sinking:5}
是的——木头是会\omterm{def:g1:floating-and-sinking:words}{漂浮}的材料，物体通常
与它的材料表现一致，如树枝或树干。
\cref{ex:g1:floating-and-sinking:log}中的原木顺河漂流。
\end{solution}

\begin{solution}{exo:g1:floating-and-sinking:6}
按压时，你会感到水把球往上推，
越深推得越用力。松手后，球会弹回
水面——水的推力占了上风。
\end{solution}

\begin{solution}{exo:g1:floating-and-sinking:7}
空心碗形：船体宽、船舷高，给水的向上推力
提供了更大的面积。\cref{ex:g1:floating-and-sinking:boat}中的橡皮泥船
就说明了这一点——材料相同，
但碗形能\omterm{def:g1:floating-and-sinking:words}{漂浮}。
\end{solution}

\begin{solution}{exo:g1:floating-and-sinking:8}
装满并盖紧的瓶子在水中\omterm{def:g1:floating-and-sinking:words}{漂浮}得较低（大多数瓶子
只是刚好浮起）；空瓶盖紧后\omterm{def:g1:floating-and-sinking:words}{漂浮}得较高，
浮在水面上——瓶内主要是空气。
\end{solution}

\begin{solution}{exo:g1:floating-and-sinking:9}
水从不会累。把橡皮泥捏成球时，
水只能向上推它底部一小块圆形面积；
做成空心船时，同样一团橡皮泥铺得很宽，
水向上推在整个底部的力足以托住它。
变的是形状，而不是水。
\end{solution}
